% Transcription — fonds Alexandre Grothendieck, shelfmark 148, batch 3 (pages 41-60).
% Established from the facsimile archives/batches/148/batch-03.pdf, cut from a
% copy of 148.pdf fetched through the project's relay
% (https://grothendieck-relay.onrender.com/source/148.pdf); PDF page k is the
% archivists' page 40+k, no cover sheet. Per the skill transcribe-grothendieck.
% The folder is seventy-one pages in four batches; this is the third. Batch 4
% (pages 61-71) was written in parallel and read for alignment (notation,
% his pagination); batches 1-2 are not yet transcribed.
%
% Pass: Opus 5.5 (claude-opus-5-5), 2026-09-26 — first pass, unchecked against the pages by a human.
%
% Pages skipped: none. Page 44 carries only two struck lines of his (the
% first draft of the opening sentence of his p. 8, page 51) and the rest of
% the sheet is blank; the two lines are given, struck. Pages summarised:
% none (no typescript). Pages 41-42 are on computer-listing paper (the
% printed listing shows through on 42; nothing of it is recorded).
%
% His own pagination: pages 41 and 42 carry no number of his. Page 43 opens
% a run he numbers 1 to 17 at top centre: his 1 = page 43, his 2 = page 45
% (page 44, the struck draft, is unnumbered), his 3) ... 17 = pages 46-60
% (he writes « 3) », « 4) », « 5) » with a parenthesis, then bare numerals;
% on page 50 the « 7 » is written over another figure). The run continues in
% batch 4 (pages 61-64 = his 18-21). Numbered runs inside it: on page 41
% a « Catalogue opérations + structures » 1)-7); on page 51 « Opérations
% essentielles », A) with (1) (circled) Gommage de repères, then on pages
% 53-54 2) Gommage de cercles de découpage and 3) Bouchage de trous; (B)
% Opération de marquage (page 55, written over a struck, boxed 4) or A));
% C) Découpage (page 57); D) Opérations de recollement (page 58, a) and b));
% and on page 60 (his p. 17) the « Récapitulation des opérations
% élémentaires » (0°), 1°)-7°), with (4°) gommage de cercle de découpage
% and (5°) bouchage de trou circled. Note for batch 4: at his p. 20
% (page 63) the gommage of cutting circles is called 5°, whereas page 60
% numbers it 4° and gives 5° to the bouchage of holes; recorded in the note
% opening page 60. No dates in his hand anywhere in the batch; \dating is
% the catalogue's « 1983 », verbatim, with no group sentence.
%
% What the batch is about. Page 41 (listing paper): a first catalogue of
% « opérations + structures » on the Teichmüller groupoids — sum, (partial)
% gommage of marking, marking, (partial) cutting, amalgamation/recollement,
% bouchage of holes (= gommage of the joining circles), gommage of inner
% circles — with the vocabulary (cercles de raccord / de découpage, sommets
% de marquage, marked and unmarked circles, an « anabelian » or « stability »
% condition, « twist » isomorphisms) and two diamonds of groupoids
% T, ST, DT, SDT -> L, SL, dL, SdL (« types combinatoires »). Page 42:
% « Les douze avatars principaux de la Tour de Teichmüller », a tower
% T_top <- T_diff <- T_an <- T_hyp <- T_C-alg <- T_Qbar-alg <- T_comb <-
% T_Lego <- T_P (topological, differentiable, C-analytic, hyperbolic,
% C-algebraic, Qbar-algebraic, combinatorial via oriented maps, « Légo »,
% « pantalonnique » with only triads as pieces), with primed variants
% T'_diff, T'_an, T'_hyp (marked points and tangent vectors). Pages 43-60
% (his 1-17): « Groupoïdes de Teichmüller » — T^+_top = T_top (compact
% oriented surfaces with boundary, isotopy classes of isomorphisms) and
% T^±_top with the signature functor to [±1]; the « most complex variant »
% S_*TD^+_top, objects (X, Σ~, R) with Σ~ = ∂X ⊔ Σ a compact 1-submanifold
% and R ⊂ Σ~ finite (joining circles, cutting circles, rigidification
% vertices), its full subgroupoids T_top, S_*T_top, TD_top; the functor
% Γ to bi-weighted (Mumford-Deligne) graphs (S, A→, σ, o, γ, r), which
% induces a bijection on π_0 and a surjection on π_1; the sum operation;
% π_0(X) ≃ π_0(Γ) and the genus formula g = Σ g_s + rg H^1(Γ) (margin);
% the special cases S_*T^±_top (isolated vertices with free edges) and
% T^±_top (a pair (g, ν) in the connected case); the refinement by
% torsors R_a under Z/r(a) ∧ ω(a) and the « graphe pondéré à repères »
% (R→ with a D_∞-action over A→) and the functor to (Gr.pondrep); then
% the essential operations — gommage of markers, of cutting circles
% (contraction of edges: a « generalisation » morphism of graphs),
% bouchage of holes (gommage of free edges), mixed gommage, surmarquage
% (with uniform gommage partiel / surmarquage as restriction / extension of
% the groups Z_a/r(a)), découpage along cutting circles (cutting an edge
% into two free edges), and recollement (an involution on free edges plus
% orientation-reversing isomorphisms, needing A' ⊂ A^+, r(a) = r(a') and an
% element of R_a ∧ R_a'); page 60 recapitulates them as 0°-7° with their
% (ir)reversibility in the margin.
%
% Notation fixed once for the batch, aligned with batch 4: his script T as
% \mathcal{T}; the prefixes S_* (« spécial », the star under the S as he
% writes it here) and the suffix D (« divisé »), top as \mathrm{top}; ± and +
% as superscripts; [±1] for the one-object groupoid; Σ~ as
% \widetilde{\Sigma}; his arrowed sets as \vec{A}, \vec{R}; the natural
% numbers as \mathbf{N}; Z_a for Z ∧_{±1} ω(a). On pages 41-42 the T is
% drawn with a doubled stem and is set \mathbb{T} (the letter of the second
% diamond on p. 41, read \mathbb{L}, is doubtful). The graph functor is
% Γ or Γ_MD as he writes it.
%
% Hardest pages: 52 (a long oblique marginal NB filling the left margin,
% mostly illegible), 46 (a block cancelled by a large cross, and two margin
% notes), 55 (heavily reworked interlinear text), 42 (the listing paper and
% the struck parentheses), 60 (margin annotations). Readings to check
% first: the first letter of « RTD... » and the « 10 opérations » on
% p. 42; the oblique margin notes of pp. 43, 48, 52, 55; « renforcement »
% and « pas un affaiblissement » on p. 56; the long sentence at the foot
% of p. 57; « proposons en vue » on p. 58; the words after « réversible à »
% in the margins of p. 60.
%
\documentclass[11pt,a4paper]{article}
% Le préambule commun à toutes les transcriptions du fonds Grothendieck.
%
% Il tient en une idée : **l'incertitude se compose, elle ne se résout pas.**
% Une transcription qui lisse un mot douteux a détruit la seule chose pour
% laquelle on la faisait — un lecteur doit pouvoir savoir, à chaque endroit,
% ce qui était sur le feuillet et ce qui a été ajouté. Les six macros qui
% suivent sont donc l'appareil critique, et non de la décoration.
%
% Les mêmes commandes sont reconnues par `scripts/render.mjs`, qui produit la
% vue de lecture du site. Une macro ajoutée ici doit l'être aussi là-bas,
% faute de quoi le rendu échouera bruyamment — ce qui est voulu : un rendu
% silencieusement faux est pire qu'un rendu absent.

\ProvidesPackage{grothendieck}[2026/08/08 Transcriptions du fonds Grothendieck]

\usepackage{amsmath,amssymb,amsthm}
\usepackage{mathrsfs}
\usepackage{stmaryrd}
% Les diagrammes commutatifs sont partout dans ce fonds : c'est la langue
% même de la géométrie algébrique de Grothendieck, et une transcription qui
% les rendrait en prose ne transcrirait rien.
\usepackage{tikz-cd}
% Français par défaut — c'est la langue des feuillets — mais l'anglais est
% chargé aussi : la traduction et le résumé sont des documents anglais, et la
% typographie française (espaces avant les deux-points) y serait une faute.
% Ces documents basculent par \selectlanguage{english} en tête de corps.
\usepackage[english,french]{babel}
\usepackage{fontspec}
\usepackage{geometry}
\geometry{a4paper,margin=25mm,marginparwidth=28mm}
\usepackage{marginnote}
\usepackage{xcolor}
% ulem, not soul. `soul`'s \st and \ul are built on a character-by-character
% scan that XeLaTeX's font machinery breaks: the document fails with an
% undefined control sequence rather than degrading. ulem does the same two jobs
% and is Unicode-engine safe. `normalem` stops it hijacking \emph, which the
% transcriptions use for Grothendieck's own emphasis.
\usepackage[normalem]{ulem}
\usepackage{hyperref}
\hypersetup{colorlinks=true,linkcolor=black,urlcolor=black,pdfborder={0 0 0}}

% --- Métadonnées du lot -----------------------------------------------------
% Reprises telles quelles par le script de rendu, qui les affiche en tête de
% la vue de lecture. Elles fixent ce que le fichier prétend couvrir : sans
% elles, une transcription flotte sans point d'attache dans un fonds de seize
% mille feuillets.

\newcommand{\folder}[1]{\gdef\grothFolder{#1}}
\newcommand{\batch}[1]{\gdef\grothBatch{#1}}
\newcommand{\leaves}[2]{\gdef\grothFirst{#1}\gdef\grothLast{#2}}
\newcommand{\foldertitle}[1]{\gdef\grothFoldertitle{#1}}
\gdef\grothFolder{?}\gdef\grothBatch{?}\gdef\grothFirst{?}\gdef\grothLast{?}\gdef\grothFoldertitle{}

% --- Le résumé --------------------------------------------------------------
%
% Il ouvre la lecture modernisée, sous le titre « Résumé », et il est composé
% en petit. Ce n'est pas un ornement : c'est le seul passage du document qui
% ne soit pas des mathématiques, et un lecteur doit pouvoir le reconnaître
% avant de l'avoir lu — puis le sauter s'il connaît déjà le sujet, ou s'y
% arrêter s'il ne le connaît pas. Le filet qui le suit marque la reprise du
% corps.

\newenvironment{resume}
  {\small\setlength{\parskip}{0.45em}}
  {\par\medskip\hrule\medskip}

% --- Mots-clés ---------------------------------------------------------------
%
% En anglais, en clôture du résumé de la lecture modernisée : le vocabulaire
% moderne sous lequel chercher ce que les feuillets construisent. Le manifeste
% du site (`npm run manifest`) les extrait pour étiqueter la cote entière —
% cette ligne est la source unique des tags, il n'y a pas de fichier à côté.
\newcommand{\keywords}[1]{\par\smallskip\noindent
  {\footnotesize\textsc{Keywords} --- \textit{#1}}\par}

% --- L'appareil critique ----------------------------------------------------

% Le numéro de feuillet, en marge. C'est l'ancre qui permet de régler un
% désaccord sur une formule en regardant *un* feuillet plutôt qu'une chemise
% de six cents. Le site s'en sert aussi pour tourner les pages du fac-similé
% quand on fait défiler la transcription.
\newcommand{\leaf}[1]{%
  \marginnote{\footnotesize\color{gray!70}\textsf{#1}}%
  \label{leaf:#1}\ignorespaces}

% Illisible : on ne devine pas. Un mot inventé qui se lit comme les autres est
% le pire résultat possible.
\newcommand{\ill}{\textcolor{red!60!black}{[\,\ldots\,]}}

% Lecture douteuse : le mot est donné, et signalé comme incertain.
\newcommand{\uncertain}[1]{\uline{#1}}

% Ajout de l'éditeur — une lettre manquante, un mot sous-entendu.
\newcommand{\add}[1]{\textcolor{blue!50!black}{[#1]}}

% Biffé par Grothendieck lui-même. Ce qu'il a rayé fait partie du document :
% une rature dit souvent où la pensée a changé de direction.
\newcommand{\struck}[1]{\sout{#1}}

% Note du transcripteur, sur l'état matériel du feuillet.
\newcommand{\note}[1]{\par\smallskip{\small\color{gray!60!black}\textit{[#1]}}\par\smallskip}

% Note marginale de Grothendieck, distincte du corps du texte.
\newcommand{\marginal}[1]{\par\smallskip\noindent\hspace{1em}%
  {\small\color{gray!60!black}#1}\par\smallskip}

% --- Titre ------------------------------------------------------------------

\AtBeginDocument{%
  \begin{center}
    {\large\bfseries Fonds Alexandre Grothendieck --- cote n\textsuperscript{o}~\grothFolder}\\[2pt]
    {\itshape\grothFoldertitle}\\[4pt]
    {\small feuillets \grothFirst--\grothLast\ (lot \grothBatch)}\\[2pt]
    {\footnotesize\color{gray!60!black}Transcription établie d'après le fac-similé numérisé
     par l'Université de Montpellier.\\
     Les lectures incertaines sont soulignées, les illisibles notées [\,\ldots\,],
     les ajouts entre crochets.}
  \end{center}
  \vspace{1em}}

\endinput


\folder{148}
\batch{3}
\pages{41}{60}
\dating{1983}
\watermark{Édition de démonstration}
\foldertitle{Teichmüller (1983) : notes manuscrites (1983)}

\begin{document}

\section{Opérations et structures}

\page{41}
\note{feuillet de papier listing, à l'encre ; les deux pages 41 et 42 ne
portent pas de numéro de sa main}

Opérations de recollage\,/\,sommes

\begin{itemize}
  \item[] objets ; diagrammes
  \item[] cercles de raccord (ou de recollage)
  \item[] cercles de découpage
  \item[] sommets de \emph{marquage} (périphériques\,/\,internes)
  \item[] cercles marqués et non marqués
\end{itemize}
\note{ces quatre dernières lignes sont réunies par une accolade à gauche}

[condition \struck{\ill{}} « anabélienne » ou de « stabilité »]

\begin{tikzcd}[column sep=small, row sep=small]
 & \mathcal{T} \arrow[dl, no head] \arrow[dr, no head] \arrow[rrrr] & & & & \mathbb{L} \arrow[dl, no head] \arrow[dr, no head] & \\
S\mathcal{T} \arrow[dr, no head] \arrow[rr, no head] & & D\mathcal{T} \arrow[dl, no head] \arrow[rr] \arrow[rrrr, bend right=20] & & S\mathbb{L} \arrow[dr, no head] & & d\mathbb{L} \arrow[dl, no head] \\
 & SD\mathcal{T} \arrow[rrrr] & & & & Sd\mathbb{L} &
\end{tikzcd}

\note{deux losanges dessinés à droite de la page, reliés par des flèches ;
le sommet du premier, $\mathcal{T}$, est précédé d'une lettre surchargée
illisible. Sous le second : « types combinatoires ». La lettre du second
losange, à double jambage, est lue $\mathbb{L}$ sans certitude ; la
minuscule $d$ de $d\mathbb{L}$, $Sd\mathbb{L}$ est telle qu'il l'écrit.
Le trait entre $S\mathcal{T}$ et $D\mathcal{T}$ n'a pas de pointe ; la
flèche courbe part de $D\mathcal{T}$ vers $d\mathbb{L}$, la flèche droite
issue de la même région aboutit près de $S\mathbb{L}$}

Catalogue opérations + structures

\begin{enumerate}
  \item[1)] Somme
  \item[2)] gommage (partiel) du marquage $\updownarrow$
  \item[3)] marquage
  \item[4)] découpage (partiel)
  \item[5)] \emph{amalgamation} ou recollage \struck{\ill{}} (sans gommage !) $\updownarrow$
  \item[6)] bouchage de trous (= gommage des cercles de raccord)
  \item[7)] gommage des cercles internes
\end{enumerate}
\note{1) à 3) sont réunis par une accolade, 4) à 6) par une autre ; les
flèches doubles verticales, à droite de 2) et de 5), semblent apparier 2)
avec 3) et 4) avec 5)}

foncteurs et iso « twist »

\page{42}
\note{papier listing, à l'encre. La lettre $\mathbb{T}$ de cette page est
tracée avec un double jambage}

Les douze avatars \struck{\ill{}}principaux de la Tour de Teichmüller :
\struck{système de groupoïdes}

groupoïde $\mathbb{R}\mathbb{T}D\ldots$\note{la première lettre, lue
$\mathbb{R}$, est douteuse} avec structures (« les 10 opérations »)
\struck{\ill{}}, ou autres\note{« ou autres » en interligne, après la
parenthèse} diverses dans ce groupoïde, dont le \struck{représentant} le plus
élémentaire est \struck{\ill{}} $\mathbb{T}_{\mathrm{top}}$

\begin{tikzcd}[column sep=small, row sep=small]
 & \mathbb{T}_{\mathrm{top}} \\
\mathbb{T}'_{\mathrm{diff}} \arrow[r, leftrightarrow] & \mathbb{T}_{\mathrm{diff}} \arrow[u] \\
\mathbb{T}'_{\mathrm{an}} \arrow[r, leftrightarrow] & \mathbb{T}_{\mathrm{an}} \arrow[u] \\
\mathbb{T}'_{\mathrm{hyp}} \arrow[u] \arrow[r, leftrightarrow] & \mathbb{T}_{\mathrm{hyp}} \arrow[u] \\
 & \mathbb{T}_{\mathbf{C}\text{-alg}} \arrow[u] \\
 & \mathbb{T}_{\overline{\mathbf{Q}}\text{-alg}} \arrow[u] \\
 & \mathbb{T}_{\mathrm{comb}} \arrow[u] \\
 & \mathbb{T}_{\mathrm{Lego}} \arrow[u] \\
 & \mathbb{T}_{\mathbb{P}} \arrow[u]
\end{tikzcd}
\note{colonne dessinée à gauche de la page, chaque flèche verticale
montant vers le cas précédent ; les doubles flèches entre $\mathbb{T}'$ et
$\mathbb{T}$ sont deux demi-flèches opposées. En face de chaque
$\mathbb{T}$, une ligne de texte, donnée ci-dessous dans l'ordre ; en marge
droite, « dimension » pour les quatre premières}

Les lignes, en entier :
\begin{itemize}
  \item cas topologiques (à bord \struck{\ill{}} topologiques)
    \marginal{dimension top.} ;
  \item cas différentiable (à \struck{pts \ill{}} bord lisse\add{s} ?)
    \struck{(ou \ill{} des pts repères)}\marginal{dim.\ diff} ;
    à gauche, $\mathbb{T}'_{\mathrm{diff}}$ : « à pts marqués et vect.\
    tangents correspondants » ;
  \item cas $\mathbf{C}$-analytiques (à bord \uncertain{analytiques})
    \marginal{dim.\ $\mathbf{C}$-\ill{}} ;
  \item cas hyperbolique (à bord géodésiques)\marginal{dim.\ hyp.} ;
  \item cas $\mathbf{C}$-algébrique (courbes MD, munies d'une
    \struck{rigidification} \ill{} des propagations de \uncertain{secondes}
    aux pts doubles, plus pts marqués) ;
  \item cas $\overline{\mathbf{Q}}$-alg.\ (propagations
    $\underline{\mathcal{O}}$ \uncertain{descendant}
    \struck{\ill{}}$/\overline{\mathbf{Q}}$) ;
  \item cas combinatoire, via cartes orientées \emph{à bord} \ill{}, avec
    sommets marqués ;
  \item cas combinatoire « Légo », avec \uncertain{10} pièces du jeu ;
  \item cas combinatoire « pantalonnique », avec comme seules pièces les
    \struck{pantalons} triades.
\end{itemize}

\section{Groupoïdes de Teichmüller}

\page{43}
\note{p.~1 de l'auteur. Ici commence une suite paginée par lui de 1 à 17,
en chiffres arabes en haut au centre, sur feuillets blancs : sa p.~1 est la
page 43, sa p.~2 la page 45, puis ses pp.~3 à 17 les pages 46 à 60 ; la
suite continue dans le lot 4 (pp.~61 à 64 = ses pp.~18 à 21). Le titre de
section est le sien}

\emph{Groupoïdes de Teichmüller.}

(1)\note{le 1 est cerclé} Cas topologique.

La variante la plus simple est
\[
\mathcal{T}^{+}_{\mathrm{top}} \quad\text{ou}\quad \mathcal{T}_{\mathrm{top}}
\]
(\emph{groupoïde de Teichmüller orienté topologique}) :

\emph{objets} : surfaces \emph{à bord} compactes orientées (\ill{}),

\emph{flèches} : classes d'isotopie d'iso (conservant l'orientation).

Variante immédiate
\[
\mathcal{T}^{\pm}_{\mathrm{top}}
\]
(\emph{groupoïde de Teichmüller biorienté topologique}) : mêmes objets,
flèches sont les classes d'isotopie d'homéomorphismes (ne conservant pas
néc.\ l'orientation).

Soit $[\pm 1]$ le groupoïde ayant un seul objet $e$ (par exemple
$e = \emptyset$) et comme groupe d'automorphismes $\pm 1$. On a un
morphisme de signature
\[
\mathcal{T}^{\pm}_{\mathrm{top}} \xrightarrow{\ \chi\ } [\pm 1]
\]
dont la \emph{fibre} \struck{est} (en \uncertain{n'importe}\ \ill{}) est
$\mathcal{T}^{+}_{\mathrm{top}} = \mathcal{T}_{\mathrm{top}}$.%
\marginal{\uncertain{canulé}, si on ne se borne \uncertain{à des surfaces
connexes} !}\note{note de la marge gauche, écrite en oblique en face de la
ligne de la fibre ; la lettre au-dessus de la flèche est lue $\chi$}

La variante la plus complexe est
\[
S_{*}\mathcal{T}D^{+}_{\mathrm{top}} \quad\text{ou}\quad S_{*}\mathcal{T}D_{\mathrm{top}},
\]
« \emph{groupoïde de Teichmüller} \add{orienté}\note{« orienté » en
interligne} \emph{spécial divisé topologique} ».

\marginal{\uncertain{Finalement}, pour \ill{} \uncertain{travaux} \ill{}
\uncertain{ultérieurs}, il semble que $\mathcal{T}^{\pm}$ \uncertain{soit}
le plus utilisé ; c'est lui qui \uncertain{mérite} donc ici la notation la
plus simple $\mathcal{T}$, et le groupoïde orienté s'appellera
$\mathcal{T}^{+}$.}\note{note écrite en oblique dans l'angle supérieur
droit de la page, en face des lignes sur $\mathcal{T}^{+}_{\mathrm{top}}$ ;
lue sur l'image tournée. Dans la suite de ces pages il écrit pourtant
encore $\mathcal{T}$ pour le groupoïde orienté et $\mathcal{T}^{\pm}$ pour
le biorienté}

\page{44}
\struck{On a ainsi la notion de graphe pondéré \emph{enrichi} (pondéré par
$\gamma$, la structure}\note{deux lignes seulement, en haut de la page,
barrées par des traits obliques ; le reste de la page est blanc. La même
phrase est reprise au début de sa p.~8 (page 51) sous la forme « graphe
pondéré à repères »}

\page{45}
\note{p.~2 de l'auteur}
\marginal{$X, \struck{\partial X}, \Sigma, S$}\note{en haut à droite}

\emph{objets} : Triples $(X, \widetilde{\Sigma}, R)$\note{$R$ écrit
au-dessus d'un $S$ biffé}, où

\quad $X$ surface à bord compacte orientée,

\quad $\widetilde{\Sigma}$ sous-variété compacte de $X$ \struck{\ill{}},
$\widetilde{\Sigma} \supset \partial X$,

\quad \struck{\ill{}} $R \subset \Sigma$ \struck{\ill{}}, partie finie.

\emph{flèches} : classes d'isotopie d'homéom.\ conservant l'orientation.

\marginal{$\widetilde{\Sigma} = \partial X \sqcup \Sigma$. On peut
considérer que $R$ rigidifie, \ill{} $\widetilde{\Sigma}$ \ldots\ le
« découpage » par $R$}\note{note oblique de la marge gauche}

\struck{Variante} Les compos.\ connexes de
$\widetilde{\Sigma} \setminus \partial X$ ($= \Sigma$) \add{et} celles de
$\partial X$ sont les \emph{cercles de raccord}\note{la phrase est
remaniée en interligne ; l'ordre des morceaux est incertain}, les
\emph{cercles de découpage}, les éléments de $R$ sont les
\emph{sommets de rigidification}\note{« rigidification » écrit au-dessus
de « marquage » biffé ; en interligne aussi, barré, « (ou de
rigidification) »} \struck{, les compos.\ connexes de $\partial X$ sont les
« trous »}. Les composantes connexes de $X$ sont les \emph{composantes} de
l'objet envisagé.

Variante $S_{*}\mathcal{T}D^{\pm}_{\mathrm{top}}$, avec

\begin{tikzcd}
S_{*}\mathcal{T}D_{\mathrm{top}} \arrow[r, hook] & S_{*}\mathcal{T}D^{\pm}_{\mathrm{top}} \arrow[r] & {[\pm 1]}
\end{tikzcd}

\ill{} \uncertain{fibre}.

On a des sous-groupoïdes pleins
\begin{itemize}
  \item[] $R = \emptyset$, $\Sigma = \emptyset$ \struck{\ill{}} :
    $\mathcal{T}_{\mathrm{top}}$ et $\mathcal{T}^{\pm}_{\mathrm{top}}$ ;
  \item[] $\Sigma = \emptyset$ \struck{\ill{}}, $R$ quelconque :
    $S_{*}\mathcal{T}_{\mathrm{top}}$ et
    $S_{*}\mathcal{T}^{\pm}_{\mathrm{top}}$ (\emph{groupoïdes de Teichmüller
    spéciaux} (ou rigidifiés) (topologiques orientés)) ;
  \item[] $R = \emptyset$ : $\mathcal{T}D_{\mathrm{top}}$ et
    $\mathcal{T}D^{\pm}_{\mathrm{top}}$ (Cas pas intéressant je crois, sauf
    pour $\Sigma = \emptyset$ également).
\end{itemize}

\marginal{MD = « modulaire » Mumford--Deligne}\note{oblique, en marge
gauche en bas de la page}

\emph{Détermination} du $\pi_0(S_{*}\mathcal{T}D_{\mathrm{top}})
\xrightarrow{\sim} \pi_0(S_{*}\mathcal{T}D^{\pm}_{\mathrm{top}})$.
On introduit les \emph{graphes} \struck{MD} bipondérés

\page{46}
\note{p.~3) de l'auteur}
qui forment un groupoïde \struck{\ill{}} (Gr.bip) ou (Gra.bipond)
\add{+ structure suppl.}\note{en interligne, comme « finis » sous
« groupoïde »}

\emph{Objets} : graphes \add{finis} $(S, \vec{A}, \sigma, o)$ (permettant
aussi des sommets isolés, des boucles fermées, + plusieurs arêtes entre deux
sommets, et aussi des « \emph{arêtes libres} » correspondant aux pts fixes
de $\sigma$) (« sommets », « arcs ») : $S$, $\vec{A}$ ens.\ finis,
$\sigma : \vec{A} \to \vec{A}$ involution, $o : \vec{A} \to S$ appl.\
« origine ». On se donne de plus
\[
S \xrightarrow{\ \gamma\ } \mathbf{N}
\]
(application « genre » ou « poids »)\note{le $\gamma$ surcharge une lettre
biffée}
\[
\vec{A}^{\sigma} = A_{\ell} \xrightarrow{\ r\ } \mathbf{N}
\]
(application \struck{\ill{}} « indice de rigidification »)\note{à gauche de
$A_{\ell}$, une lettre hachurée et $\overset{\mathrm{def}}{=}$ ; la lecture
$\vec{A}^{\sigma}$ est de l'édition, d'après la suite de la page}

\note{ici un bloc de lignes est barré en croix par de grands traits
obliques : $\struck{\ldots\,\vec{A}/\sigma}$ « ens.\ des sommets
supplémentaires » ; « \ill{} encore, posant
$\widetilde{S} = S \sqcup A_{\ell}$ (identifié \ldots\ l'ens.\ des
\ldots\ correspondant aux $a \in A_{\ell}$) \ldots\ les « sommets »,
\struck{« marqués » et} « non marqués » (i.e.\ $\in S$) \ldots\ libres, du
graphe), \ldots\ on donne $\widetilde{S} \xrightarrow{\ \pi\ } \mathbf{N}$
(\emph{pondération} du graphe) »}

\marginal{pas assez précis, comme donnée : la pondération d'une arête doit
être remplacée par celle d'un pol.\ comb.\ $T_a$ avec
$\omega(T_a) \simeq \omega(a)$}\note{note oblique de la marge gauche ; la
fin est lue sans certitude}

\marginal{j'ai indiqué une pondération des sommets, mais pas des arêtes}
\note{note encerclée, marge gauche}

\note{dessin en marge gauche : un graphe à sommets pondérés $0$, $13$, $4$,
$2$, $0$, $3$, $1$, avec arêtes multiples, une boucle en $0$ et des arêtes
libres ; au-dessous, sous « 2° », un sommet $0$ avec deux arêtes libres, et
un sommet $3$ portant une boucle}

\begin{align*}
\widetilde{S} &= S \sqcup A_{\ell}, \\
\widetilde{A} &= \vec{A}/\sigma = A \sqcup A_{\ell}, \\
\vec{A} &= \vec{A}_0 \sqcup \vec{A}_{\ell},
  \qquad \vec{A}_0 = \vec{A} \setminus \vec{A}^{\sigma},
  \quad \vec{A}_{\ell} = \vec{A}^{\sigma},
\end{align*}
\note{$A_{\ell}$ est écrit sous $S \sqcup$, relié par un signe d'inclusion vertical à un
symbole biffé}

\begin{tikzcd}
\vec{A} \arrow[d] & = & \vec{A}_0 \arrow[d, "\deg 2"] & \sqcup & \vec{A}_{\ell} \arrow[d, "\deg 1"] \\
\widetilde{A} & = & A & \sqcup & A_{\ell}
\end{tikzcd}

\note{la flèche verticale de gauche est doublée d'un faisceau de traits
(une flèche hachurée)}

Ordre d'un sommet
\[
\sigma(s) = \sigma_{\ell}(s) + \sigma_0(s),
\]
$\sigma_{\ell}$ correspondant aux « trous » ou cercles de recollage.

La donnée des \emph{graphes} sous-jacents à un graphe \struck{MD} bip.\
revient à celle des graphes « ordinaires » correspondants (sans arêtes
libres, i.e.\ où $\sigma$ opère sans pts fixes), i.e.
\[
\bigl(S,\ (\vec{A} \setminus \vec{A}^{\sigma}),\
\sigma\,|\,(\vec{A} \setminus \vec{A}^{\sigma}),\
o\,|\,(\vec{A} \setminus \vec{A}^{\sigma})\bigr),
\]
\emph{plus} un ensemble \add{fini} $A_{\ell}$ (arêtes \struck{\ill{}} libres)
\struck{ou « sommets-trous »} et une application
\[
A_{\ell} \longrightarrow S .
\]

\emph{Flèches} : Les iso.\ de telles structures.

\page{47}
\note{p.~4) de l'auteur}
On a un foncteur
\[
(*)\qquad \Gamma : S_{*}\mathcal{T}D^{\pm}_{\mathrm{top}} \longrightarrow
\text{(Gra bipond)}
\]
\note{avant « (Gra bipond) », un sigle biffé, peut-être « FMD »}
qui associe à un \emph{objet} $(X, \widetilde{\Sigma}, R)$ de
$S_{*}\mathcal{T}D^{\pm}_{\mathrm{top}}$ le graphe \add{MD}
$(S, \vec{A}, \sigma, \struck{\ill{}}\, o, \rho)$\note{« MD » en
interligne ; sur $o$, une lettre surchargée et hachurée ; au-dessus,
$\sigma$} suivant :
\begin{itemize}
  \item[] $S = \pi_0(X \setminus \widetilde{\Sigma}) \rightleftarrows
    \pi_0(\mathrm{Déc}(X, \widetilde{\Sigma}))$ ;
  \item[] $\vec{A}$ = ens.\ des couples $(C, \omega)$, où
    $C$ \struck{est \ill{}} $\in \pi_0(\widetilde{\Sigma})$, et
    $\omega \in \mathrm{Or}(C) \simeq \mathrm{rive}(C)$\note{au-dessus de
    $\simeq$ : « rive$(C)$ » et un $x$, avec une accolade renvoyant à
    « sommets »}, \ill{} : la condition que si $C \in \pi_0(\partial X)$,
    donc on soit l'orientation induite, i.e.\ la rive qui se trouve dans
    $X$ ;
  \item[] $\sigma$ = échange des rives (quand on a le droit d'échanger,
    i.e.\ si $C \notin \pi_0(\partial X)$, \ill{} $\sigma(a) = a$) ;
  \item[] $o$ associe à $(C, \text{rive de } C)$ l'unique « face »
    $\in \pi_0(X \setminus \widetilde{\Sigma})$ qui contient cette rive ;
  \item[] $\gamma : S \to \mathbf{N}$ (où $S = \pi_0(X \setminus
    \widetilde{\Sigma})$) = application \emph{genre} des composantes $X_i$ ;
  \item[] $r : A = \pi_0(\widetilde{\Sigma}) \to \mathbf{N}$ :
    $C \mapsto \operatorname{Card}(R \cap C)$.
\end{itemize}
\note{la lettre notée $\rho$ dans le sextuple devient $r$ à la dernière
ligne ; le $\gamma$ est écrit au-dessus d'une flèche courbe}

Le foncteur $(*)$ \struck{est surjectif et} induit une \emph{iso sur les
$\pi_0$} et une \emph{surjection} sur les $\pi_1$\note{« induit une iso sur
les » en interligne, au-dessus du passage biffé} (c'est une
\uncertain{$0$}-isomorphisme, en langage homotopique).

\struck{\ill{}}

\emph{Remarques.} Dans $S_{*}\mathcal{T}D^{\pm}_{\mathrm{top}}$ il y a une
\emph{opération} évidente de \emph{somme}, dans (Gra bipond)
\struck{FMD} aussi, et $(*)$ y commute.

\page{48}
\note{p.~5) de l'auteur}
Cette opération de somme est la plus triviale parmi les opérations qu'il
nous faudra expliciter et transcrire (combinatoirement) dans la suite.
Notons de suite, pour $(X, \widetilde{\Sigma}, R)$ donné :
\[
X \ \text{connexe} \iff \Gamma_{\mathrm{MD}}(X, \widetilde{\Sigma}, R)\ \text{connexe}
\]
\note{après $X$, une lettre biffée ; l'indice de $\Gamma$, lu MD, est
surchargé}
plus précisément
\[
\pi_0(X) \simeq \pi_0(\Gamma_{\mathrm{MD}}(X, \widetilde{\Sigma}, R)).
\]

\marginal{\ill{} plus généralement, donnons la formule pour le genre d'une
\uncertain{composante} $\alpha \in \pi_0(X) \simeq \pi_0(\Gamma)$ en termes
de la pondération : \emph{le genre total} est
$g = \sum_{s \in S} g_s + c(\Gamma)$, $c(\Gamma)$ = l'ordre de connexion du
graphe \struck{\ill{}} $= \operatorname{rg} H^1(\Gamma)$}\note{longue note
oblique de la marge gauche, lue sur l'image tournée ; l'ordre des
morceaux est celui que suggère le tracé, et la fin
« $= \operatorname{rg} H^1(\Gamma)$ » est écrite au-dessus d'un mot
biffé}

\emph{Cas particuliers.}

a) Cas de $S_{*}\mathcal{T}^{\pm}_{\mathrm{top}}$ \struck{\ill{}} --- i.e.\
$\Sigma = \emptyset$, $\widetilde{\Sigma} = \partial X$.

Cela correspond aux graphes MD qui n'ont que des arêtes libres, i.e.\ tous
les sommets sont isolés (avec \add{des} arêtes libres) : ce graphe est donc
déterminé par
\struck{\ill{}}

\begin{tikzcd}
A \ (= \vec{A} \setminus \vec{A}^{\sigma}) \arrow[d] \arrow[r, "r"] & \mathbf{N} \\
S \arrow[r, "\gamma"] & \mathbf{N}
\end{tikzcd}

\note{dessin en marge gauche : des sommets isolés pondérés $g_0$, $g_1$,
$g_2$, $g_3$, $g_4$, portant chacun quelques arêtes libres, réunis par une
accolade. La parenthèse sous $A$ se lit
$\vec{A} \setminus \vec{A}^{\sigma}$ sans certitude}

Le cas de $\mathcal{T}^{\pm}_{\mathrm{top}}$ correspond : de plus
$\rho = 0$, donc la situation est déterminée par

\begin{tikzcd}
A = A_0 \arrow[d] & \\
S \arrow[r, "\gamma"] & \mathbf{N}
\end{tikzcd}

Le cas connexe ($\operatorname{card} S = 1$, i.e.\ $S = \lbrace s \rbrace$)
\struck{type $(g, \nu)$ d'une structure \ill{}}\note{deux lignes biffées}
est donné par la donnée de deux entiers $(g, \nu)$ ($g$ = genre, $\nu$ = nb
de trous) :
\[
g = \gamma(s), \qquad \nu = \operatorname{card} A .
\]

\page{49}
\note{p.~6 de l'auteur}
\emph{Remarques.} En fait, on a un foncteur plus précis, vers une catégorie
de graphes munis d'une structure où la partie $\rho$ de la pondération est
\struck{\ill{}} \add{\ill{}}\note{un mot biffé, un autre en interligne,
tous deux illisibles} par une structure plus précise, à savoir
\struck{une \ill{} application} \add{une famille}

\[
\struck{A\,(\widetilde{R}, \alpha,\ \rho : R \to \vec{A})}\qquad (R_a)_{a \in A}
\]
\note{la première formule, barrée, portait sous $\vec{A}$ un « $=\vec{A}_0$ »
biffé}

\marginal{$\mathrm{def}\ \omega(a) = \lbrace \tilde{a} \in \vec{A} \mid
\tilde{a}\ \text{au-dessus de}\ a \rbrace$}\note{oblique, marge gauche}

\struck{disjointes, et pour tout $a \in A$, si $\omega(a)$ désigne l'un
des \ill{} (card $\omega(a) \in \lbrace 1, 2 \rbrace$), une application
\ill{} $R$, et $\rho$ une application $R \to A$ \ill{}}\note{quatre lignes
barrées}

\[
\omega \mapsto u_{\omega} : \omega(a) \longrightarrow \mathrm{Circ}(R_a)
\]
\ill{} \uncertain{circulaires} sur $R_a$\note{« $\omega \mapsto u_\omega$ »
est ajouté en interligne ; au-dessus de la ligne biffée qui précède,
$\struck{R/\alpha \to \vec{A}}$}

avec la condition que si $\omega, \omega' \in \omega(a)$, $\omega \neq
\omega'$, alors $u_{\omega'} = u_{\omega}^{-1}$
\struck{$u_{\omega} u_{\omega'} = \mathrm{id}_{R}$}.

On peut aussi considérer que la donnée des structures revient à la donnée de
\[
r : A \to \mathbf{N}, \qquad a \mapsto \operatorname{card} \rho^{-1}(a)
\]
(la « pondération-arêtes »), et à la donnée d'une famille de torseurs
\[
(R_a)_{a \in A},
\]
$R_a$ un torseur sous
$\mathbf{Z}/r(a)\mathbf{Z} \wedge^{\omega(a)}_{\lbrace \pm 1 \rbrace}$.
\note{le $\omega(a)$ est écrit en exposant du $\wedge$, $\lbrace \pm1
\rbrace$ en indice}

On associe à $(X, \widetilde{\Sigma}, R)$ une telle structure
supplémentaire :
\[
R_a = R \cap \widetilde{\Sigma}_a
\]
($\widetilde{\Sigma}_a$ la comp.\ connexe de $\widetilde{\Sigma}$ qui
correspond à $a \in A \simeq \pi_0(\widetilde{\Sigma})$),
\struck{$R = \ldots$, $R \to S = \pi_0(\widetilde{\Sigma})$}, $\omega
\mapsto u_{\omega}$ défini de façon évidente,
\struck{et $\rho$ déduit de l'ordre circulaire dans les $R_a$}.

\emph{NB.} Même si $r(\omega) = 1$ ou $r(\omega) = 2$, le fait de disposer

\page{50}
\note{p.~7 de l'auteur}
d'une structure déjà orientée si $a$ est une arête libre (correspondant à
un $C \in \pi_0(\partial X)$), ou qui dispose de l'une des deux orientations
$\omega(a)$ dans le cas où $a$ est une arête ordinaire, implique que l'on
tient au \uncertain{moins}, pour chaque $a \in A$, un polygone combinatoire
$P_a$, dont l'ens.\ des sommets est $R_a$, et l'ens.\ des arêtes peut se
décrire \struck{\ill{}} de façon canonique en termes de $R_a$, $\omega_a$ et
$\omega_a \to \mathrm{Circ}(R_a)$. Il peut être plus commode, \ill{}
\ill{} \uncertain{d'oublier} dans la structure, de \ill{} \ill{} \ill{}
la famille des $(R_a)$ serait remplacée par un \struck{qui serait dans un
contour} ensemble de \emph{contour combinatoire}, i.e.\ un ensemble \ill{}
\uncertain{repères}
\[
\vec{R}
\]
(plutôt « repères de rigidification »), \ill{} d'une opération de
\[
D_{\infty} = \lbrace \sigma_0, \sigma_1 \mid \sigma_0^2 = \sigma_1^2 = 1 \rbrace
\]
dessus, et une application
\[
\vec{R} \longrightarrow \vec{A}
\]
compatible avec l'hom.
\[
D_{\infty} \xrightarrow{\ \varepsilon\ } \lbrace \pm 1 \rbrace
\]
sur les groupes d'opérateurs, mais de plus un sous-ens.
\[
\vec{R}^{+}_{\ell} \subset \vec{R}\,|\,\vec{A}_{\ell}
\]
($\vec{A}_{\ell}$ = arêtes libres, correspondant aux cercles de
recollement, i.e.\ les $\in \pi_0(\partial X)$) définissant une orientation
des contours correspondants.
\note{$\vec{R}$ est souligné (d'un trait ondulé) à chaque occurrence}

\page{51}
\note{p.~8 de l'auteur}
On a ainsi la notion de \emph{graphe pondéré} \emph{à repères}
\struck{enrichi} $\Gamma$ :\note{« à repères » en interligne, au-dessus
d'un mot biffé ; cf.\ les deux lignes barrées de la page 44}

\begin{tikzcd}
\vec{R}^{+}_{\ell} \arrow[r, hook] & \vec{R} \arrow[r] & \vec{A} \arrow[d, "o"'] \\
 & & S \arrow[r, "\gamma"] & \mathbf{N}
\end{tikzcd}

\note{sous $\vec{R}$ : « $\sigma_0, \sigma_1$ » ; sur $\vec{A}$, une
flèche en boucle marquée $\sigma$ ; entre $\vec{R}^{+}_{\ell}$ et $\subset$,
un « $D_{\infty}$ » biffé}

Groupoïde (\emph{Gr.pondrep}), foncteur
\[
S_{*}\mathcal{T}D^{\pm}_{\mathrm{top}} \longrightarrow (\mathrm{Gr.pondrep})
\]
qui sera important sans doute dans l'explicitation de l'opération de
« recollement ». Il est clair en tous cas que toute description
algébrico-combinatoire de $S\mathcal{T}D^{\pm}_{\mathrm{top}}$ devra
permettre, pour chaque objet de la \struck{catégorie} \add{groupoïde}
paradigme, de reconstituer (à isom.\ unique près) les \emph{graphes pondérés
à repères} correspondants --- plus précisément, il faut reconstituer le
foncteur canonique
\[
S_{*}\mathcal{T}D^{\pm}_{\mathrm{top}} \longrightarrow (\mathrm{Gr.pondrep.})
\]
(et non seulement le foncteur à valeurs dans (Gr.bipond)).

\subsection*{\emph{Opérations essentielles}}

A) Opérations « d'oubli » ou de « \emph{gommage} ».

(1)\note{cerclé} \emph{Gommage de repères.} À un objet
$(X, \Sigma, R)$, muni de plus d'une partie $R' \subset R$, il associe
$(X, \Sigma, R \setminus R')$ (gommage de $R'$).

\page{52}
\note{p.~9 de l'auteur}
\emph{NB.} Ayant (via un groupoïde équivalent de description purement
combinatoire) « reconstruit » \struck{ou plutôt} réalisé le groupoïde
$S\mathcal{T}D^{\pm}_{\mathrm{top}}$, et le foncteur
\[
S\mathcal{T}D^{\pm}_{\mathrm{top}} \longrightarrow (\mathrm{Gr.pond.rep}),
\]
\struck{\ill{}} on devra pouvoir donc aussi \add{l'équivalent de}
\struck{\ill{}} la donnée d'un $(X, \Sigma, R)$, muni \emph{de plus} d'une
partie $R' \subset R$ --- \struck{\ill{}} i.e.\ arriver à la construction
d'un groupoïde

\begin{tikzcd}
S_{*}\mathcal{T}D^{\pm}_{\mathrm{top}}\ (\text{à sous-ens.\ marqué $R'$ de sommets-repères}) \arrow[d] \\
S_{*}\mathcal{T}D^{\pm}_{\mathrm{top}}
\end{tikzcd}

il nous faudra donc \add{décrire} un foncteur « gommage de $R'$ »

\begin{tikzcd}
S_{*}\mathcal{T}D^{\pm}_{\mathrm{top}}(R'\ \text{marqué}) \arrow[r, "\text{gommage de } R'"] & S_{*}\mathcal{T}D^{\pm}_{\mathrm{top}}
\end{tikzcd}

\struck{Cela s'exprime, sur le} donnant lieu : un diagramme commutatif de
foncteurs

\begin{tikzcd}[column sep=small]
S_{*}\mathcal{T}D^{\pm}_{\mathrm{top}}(\ldots R' \ldots) \arrow[r, "\text{gommage de } R'"] \arrow[d, "\text{oubli du marquage } R'"'] & S_{*}\mathcal{T}D^{\pm}_{\mathrm{top}} \arrow[d] \\
S_{*}\mathcal{T}D^{\pm}_{\mathrm{top}} \arrow[r] & (\mathrm{Gra.bipond})
\end{tikzcd}

\note{sur la page, les deux flèches aboutissant à (Gra bipond) sont
obliques ; à droite de (Gra bipond), une flèche vers « graphes avec
pondération par $\gamma$, mais non par $r$, » ; dans (Gra bipond) une
lettre est surchargée}

Alors qu'au niveau\note{la phrase s'interrompt en bas de page ; sa p.~10
(page 53) reprend « des graphes pondérés à repères »}

\marginal{\emph{NB.} Cette opération \ill{} \ill{} d'un \emph{affaiblissement}
de structure --- mais \ill{} \ill{}, \ill{} \ill{} \ill{} donnée de
$\Sigma^{+}$, lorsqu'on « gomme » \ill{} les repères \ill{} $\Sigma$ ---
c'est au contraire plutôt dans le sens d'un \emph{renforcement} de
structure --- ou plutôt, \ill{} \ill{} \ill{} $R'_C = R' \cap C$, dans le
cas \ill{} \ill{} \ill{} \ill{} catégories \ill{} \ill{} \ill{} \ill{}.
Il faudrait donc distinguer \ill{} $S_{*}\mathcal{T}D^{\pm}$ \ill{}}%
\note{longue note oblique occupant toute la marge gauche de la page, en
écriture serrée ; seuls quelques morceaux se lisent, sur l'image tournée}

\page{53}
\note{p.~10 de l'auteur}
des graphes pondérés à repères, l'opération \struck{s'explicite} se reflète
par une opération correspondante au niveau de $\vec{R}$, par le « gommage »
d'une partie $\vec{R}'$ stable par $D_{\infty}$, i.e.\ par $\sigma_0,
\sigma_1$.

2) \emph{Gommage de cercles de découpage.}

On se donne maintenant une \struck{ens.\ de composantes connexes de}
\add{(ouv.\ et) fermée} 1-sous-variété $\Sigma'$ de $\Sigma \setminus
\partial X$\note{« (ouv.\ et) fermée » en interligne}, correspondant à un
ensemble \struck{\ill{}} de comp.\ connexes de $\Sigma'$, i.e.\ à une partie
$A'_0$ de $\pi_0(\Sigma') = A_0$ (arêtes ordinaires du graphe)\note{ainsi
sur la page : $\pi_0(\Sigma')$, là où l'on attendrait $\pi_0(\Sigma)$} ---
et on remplace $\Sigma$ par $\Sigma \setminus \Sigma'$. Au niveau des
graphes pondérés, cela se traduit par une opération de \emph{contraction en
un point} des arêtes correspondantes, d'où formation d'un nouveau graphe, et
d'un \emph{morphisme de généralisation} (à l'opposé de « spécialisation »)
de l'ancien vers le nouveau. Il convient ici d'expliciter comment la
pondération suit\ldots\ Au niveau des sommets \add{(repères)}, il est
entendu que les sommets repères se trouvant sur $\Sigma'$ sont gommés en
même temps --- donc, sur les graphes pondérés à repères, on fait sauter
l'ens.\ des repères qui sont au-dessus des arêtes $\in A'_0$ !

\page{54}
\note{p.~11 de l'auteur}
3) \struck{Gommage} \add{Bouchage de trous, ou « gommage »} des cercles de
recollement.\note{« Bouchage de trous, ou « gommage » » en interligne,
au-dessus de « Gommage » biffé}

Cette opération est relative à la donnée d'une 1-sous-variété fermée de
$\partial X$, i.e.\ d'une partie $A'_{\ell}$ de $\pi_0(\partial X) \simeq
A_{\ell}$ (ens.\ des arêtes libres du graphe). Elle consiste à boucher les
trous correspondants, au moyen des cônes sur les composantes connexes du
bord de $\Sigma'_{\ell}$. Au niveau des graphes, cela consiste à gommer les
arêtes libres correspondantes, ainsi bien entendu aussi les repères de
marquage qui se trouvent au-dessus --- \struck{mais} les pondérations des
sommets restant inchangées.

\note{ici un signe de renvoi, une croix à pied}

\emph{Gommage simultané mixte}, correspondant à une partie
$R' \sqcup \Sigma'$ \struck{$A' \sqcup R'$} de $R \sqcup \Sigma$
\struck{quelconque} \struck{des} qui soit à la fois ouverte et fermée, et
telle \struck{$A \sqcup R$} que $R' \supset R \cap \Sigma'$,
\struck{ens.\ de composantes connexes de $R \sqcup \Sigma$} ou ce qui
revient au même, une partie $R' \sqcup A'$ de $R \sqcup A$, telle que
$R' \supset R\,|\,A'$.

Au niveau des graphes pondérés à repères, cette opération consiste à enlever
\struck{\ill{}} la partie corr.\ $\vec{R}'$ de $\vec{R}$, à contracter en un
point les arêtes de $A'_{\ell}$ \struck{\ill{}} (considérées comme libres que
liées), en faisant suivre la pondération.\note{l'indice de $A'$ surcharge
un mot, lu $\ell$ sans certitude ; « considérées comme libres que liées »
est douteux}

\page{55}
\note{p.~12 de l'auteur}
(B)\note{« (B) » écrit au-dessus d'un « 4) » (ou « A) ») biffé et
encadré} \emph{Opération de} \struck{\ill{}} \emph{marquage}\note{au-dessus
de la ligne, « d'un » ; dans la ligne, un mot biffé commençant par
« aff- », et en interligne « ou \ill{} », biffé aussi}

\struck{On peut la définir par la donnée, \ill{} \ill{} composantes connexes
de \ill{} \ill{} $R$ \ill{} $\widetilde{\Sigma} \setminus R$. Soit
$A^{+}_{*} \subset A_{*}$ l'ens.}\note{trois lignes barrées d'un trait et
de hachures obliques}

Soit $A^{+}$ l'ens.\ des arêtes (libres ou liées) du $\Gamma$ \add{bip.}\
telles que $r(a) > 0$, \struck{« effectivement marquées », et $B$ l'ens.\
des}, correspondant à des $C \in \pi_0(\Sigma)$ telles que
$C \cap R \neq \emptyset$, soit $B^{+}$ l'ens.\ des comp.\ connexes des
\[
\Sigma^{+} \setminus R \cap \Sigma^{+}
\]
(lesquelles sont des intervalles ouverts), un \emph{surmarquage} est défini
par la donnée d'une application
\[
B^{+} \xrightarrow{\ r'\ } \mathbf{N},
\]
et peut se décrire comme l'opération qui \add{\ill{}} sur chaque
\add{(\ill{})} \struck{arête} \add{\ill{}} de $\Sigma$, correspondant
à un $b \in B^{+}$, $r'(b)$ nouveaux sommets de
\struck{\ill{}} rigidification. Cette fois-ci, \struck{l'opération} il y a
des choix arbitraires --- néanmoins, à un objet de $S\mathcal{T}D^{+}$,
défini \add{soit $(X, \Sigma, R_1)$} --- muni bien à isom.\ unique près ---
entendu d'un ens.\ marqué $R' \subset R_1$, tel que l'opération inverse de
gommage de $R'$ redonne l'objet de départ $(X, \Sigma, R)$.\note{la fin de
la page est très remaniée en interligne ; l'ordre des morceaux est celui de
l'édition}

\marginal{\struck{arêtes} cercles \struck{effectivement marqués} marqués,
\uncertain{arêtes} sur $\Sigma$ \ill{} (\emph{NB} elles sont les
intersections des cercles \uncertain{effectivement} marqués, et
\uncertain{sont} des intervalles ouverts)}\note{note oblique de la marge
gauche, reliée par un trait à la définition de $A^{+}$}

\page{56}
\note{p.~13 de l'auteur}
On va peut-être se borner par la suite à des opérations de gommage de
repères, ou de surmarquage, de nature \add{\uncertain{assez}} spéciale, qui
se traduisent, sur les torseurs
\[
R_a \quad (a \in A)
\]
sous les
\[
\mathbf{Z}_a/(r(a)) \qquad \bigl(\mathbf{Z}_a = \mathbf{Z} \wedge_{\lbrace \pm 1 \rbrace} \omega(a)\bigr),
\]
par une opération de \emph{restriction} des groupes d'opérateurs de
$\mathbf{Z}_a/r(a)$ à $\mathbf{Z}_a/r'(a)$, où $r'(a) \mid r(a)$
(opérations de \emph{gommage partiel uniforme} (\ill{})) ou
d'\emph{extension} des groupes d'opérateurs de $\mathbf{Z}_a/r(a)$ à
$\mathbf{Z}_a/r''(a)$, où $r''(a) = d(a)\,r(a)$ est un multiple de $r(a)$
(opération de \emph{surmarquage uniforme}, \ill{} de multiplic.\ par
$d(a)$).\note{après $\mathbf{Z}_a/(r(a))$, une formule surchargée,
illisible ; « $= d(a)r(a)$ » est écrit en interligne au-dessus de
$r''(a)$}

On se rend compte \ill{} que le gommage \emph{partiel} est en fait un
\uncertain{renforcement} de structure, \uncertain{pas} un affaiblissement,
et inversement pour les surmarquages uniformes --- qui, au niveau
combinatoire, seront définis sans donnée de structure supplémentaire (autre
bien sûr que les multiplications $d(a)$).

\page{57}
\note{p.~14 de l'auteur}
C) \emph{Découpage} \struck{de} suivant une partie \add{ouv.\ et f.}\
$\Sigma' \subset \Sigma_0$ de la sous-variété formée des cercles de
découpage --- cela correspond à une partie $A'$ de $A_0$ (arêtes
\add{liées} \struck{ordinaires}), mais cette fois-ci il ne s'agit pas de
« gommer » i.e.\ de contracter en un point, mais de « découper », ce qui au
niveau des graphes revient à découper une arête par son milieu (i.e.\ à la
remplacer par deux arêtes libres).

Le marquage suit. Mais il y \struck{\ill{} \ill{}, par un \ill{}} a ici à
distinguer entre le cas \add{où $C$} d'une opération ordinaire --- celui où
$C$ est effectivement \struck{marqué} \struck{où il y a \ill{}} i.e.\
$C \cap R \neq \emptyset$, i.e.\ $C \subset \Sigma^{+}$ --- \ill{} à peu de
chose près (une ambiguïté dans un ens.\ fini, \ill{}) peut se
« inverser », à l'aide des \ill{} \ill{}) --- et celui où $C$ n'est pas
marqué. Mais il semble que dans la façon dont nous allons travailler

\page{58}
\note{p.~15 de l'auteur}
pour \struck{\ill{}} traduire le groupoïde de Teichmüller topologique, on ne
trouvera jamais que des structures où les cercles de découpage sont
effectivement marqués. \struck{\ill{} \ill{} distinguer, \ill{} pour
$S\mathcal{T}D^{\pm}_{\mathrm{top}}$ le groupoïde} \add{Néanmoins on} veut
avoir le droit de les « gommer », d'où la \ill{} ---
\add{\uncertain{marquages}} \struck{même} s'ils y sont bel et bien !

D) \emph{Opérations de recollement.} C'est la plus \struck{\ill{}}
\add{essentielle}, la moins triviale peut-être de toutes ces opérations.
Pour ce que nous \uncertain{proposons en vue}, il faut comme donnée
préliminaire
\begin{itemize}
  \item[a)] Une partie \struck{\ill{}} $A'$ de $A_{\ell}$ (arêtes libres),
    et une involution ss pts fixes (i.e.\ une relation d'équivalence à
    classes de cardinal 2) $\sigma'$ de $A'$ \struck{\ill{}} ;
  \item[b)] Pour \struck{$\lbrace a, a' \rbrace \in A'_2$} toute orbite
    $\lbrace a, a' \rbrace$ de $A'$, une isom.\ renversant l'orientation
    de $C_a$ sur $C_{a'}$ et de $C_{a'}$ sur $C_a$, inverses l'une de
    l'autre.
\end{itemize}
Pour que l'opération soit bien déterminée à isom.\ unique

\page{59}
\note{p.~16 de l'auteur}
près dans la catégorie topologique $S\mathcal{T}D^{\pm}_{\mathrm{top}}$, il
faut \struck{\ill{}} qu'on ait
\[
A' \subset A^{+},
\]
i.e.\ que les $C_a$ ($a \in A'$) soient effectivement marqués. De plus,
quitte à avoir effectué au préalable des gommages de marquage, ou des
surmarquages (on peut \uncertain{toujours} s'en tirer par surmarquages
uniformes \ill{} \ill{}), on suppose $r(a) = r(a')$ si $a' = \sigma'(a)$,
$a, a' \in A'$. Cela dit, l'opération sera bien déterminée dans la
catégorie topologique, pourvu qu'on \struck{ait} exige que le marquage soit
respecté, et que l'on se donne l'\struck{\ill{}} anti-isomorphisme des
torseurs sous $\mathbf{Z}/r\mathbf{Z}$
\[
R_a \ \text{et}\ R_{a'} \qquad (r = r(a) = r(a')),
\]
i.e.\ un élément de
\[
R_a \wedge_{\mathbf{Z}/r\mathbf{Z}} R_{a'} .
\]

\struck{\emph{NB}, \ill{}} Au niveau des graphes pondérés, l'opération
consiste à rattacher deux à deux \add{par leurs bouts} des arêtes libres ---
la pondération des sommets n'est pas changée, et la pondération des arêtes
(resp.\ la structure à repères) suit de façon \struck{\ill{}} évidente
(resp.\ \ill{} évidente).

\page{60}
\note{p.~17 de l'auteur. Les numéros 4° et 5° sont cerclés dans la marge ;
la page se continue dans le lot 4, où ses pp.~18 à 21 (pages 61 à 64)
renvoient aux opérations 1° à 7° de cette récapitulation. Ici, 4° est le
gommage des cercles de découpage et 5° le bouchage des trous ; à sa p.~20
(page 63), le gommage des cercles de découpage est appelé 5°}

\emph{Récapitulation des opérations élémentaires}

\begin{itemize}
  \item[(0°)] Somme directe\add{s}.
  \item[1°)] \emph{gommage} \add{total} \struck{\ill{}} des repères se
    trouvant sur un \add{\uncertain{les}} cercle\add{s} de recollement, ou
    \struck{\ill{} \ill{}} de découpage $\iff$ suppression d'\add{une}
    partie $R_a$ ($a \in A$) $\iff$ modification de $r$ en le
    \struck{\ill{} $R$ se trouvant sur une partie} remplaçant par $0$ en
    $a$.\marginal{fortement irréversible ; $r_a \neq 0$}
  \item[2°)] \emph{gommage partiel uniforme d'exposant $d_a \mid r_a$} :
    Restriction des \struck{scalaires} \add{groupes de $R_a$} de
    $\mathbf{Z}_a/r_a\mathbf{Z}$ à $\mathbf{Z}_a/r'_a\mathbf{Z}_a$, où
    $r'_a \mid r_a$ (pour $a$ fixé, il y a $r_a/r'_a = d_a$ choix
    possibles, \uncertain{où} on suppose $d_a > 1$, i.e.\
    $r'_a \neq r_a$).\marginal{renforcement de structure faible ;
    $r_a \geq 1$ \struck{\ill{}} ; $r'_a \neq r_a$ et $r'_a \mid r_a$,
    i.e.\ $d_a = r_a/r'_a \neq 1$}\note{la note de marge est encadrée ;
    le $\neq$ de $r'_a \neq r_a$ surcharge un autre signe}
  \item[3°)] \emph{surmarquage uniforme d'exposant $d_a > 1$}
    \struck{Extension des} --- i.e.\ extension des groupes de $R_a$ de
    $\mathbf{Z}_a/r_a\mathbf{Z}_a$ à
    $\mathbf{Z}_a/d_a r_a\mathbf{Z}_a$.\marginal{affaiblissement de
    structure faible ; $r_a \neq 0$}
  \item[(4°)] gommage de cercle de découpage $C_a$ ($a \in A_0$).
  \item[(5°)] bouchage de trou $C_a$ ($a \in A_{\ell}$).
    \marginal{fortement irréversibles}\note{accolade en marge réunissant
    (4°) et (5°)}
  \item[6°)] Découpage le long de $C_a$ ($a \in A_0$)${}^{*}$.
    \marginal{fortement irréversible si $r_a = 0$, \struck{faible}
    réversible à \ill{} près si $r_a \geq 1$}
  \item[7°)] Recollement le long de $C_a$, $C_{a'}$ ($a, a' \in A_{\ell}$,
    $a \neq a'$, $r_a = r_{a'}$), via un élément de
    $R_a \wedge R_{a'}$.\marginal{$r_a = r_{a'} \geq 1$ ; réversible à
    \ill{} près}
\end{itemize}
\note{les notes de marge sont écrites en oblique, en face de chaque
opération, et séparées par des traits}

${}^{*}$ Il faudrait peut-être séparer 6°) en deux, suivant que $r_a = 0$ ou
$r_a \neq 0$ (ou exclure le cas $r_a = 0$ ?)

\end{document}
